\documentclass[11pt]{article}
\usepackage{mathblatt}
% Probe Serie 08.10.2026: Übersichtsblatt (Einstiegsseite) Kurvenuntersuchung, von Hand gesetzt.
% Teile und Reihenfolge: Lambacher Schweizer (EP I 2, 4, 5, 6–7; EP IV 1–6; QP I 8; QP IV 3; QP V 4),
% Stoffe aus abitur/gliederung/kurvenuntersuchung.md; Rand = höchstens 3 Aufgaben in höchstens 2 der
% Jahrgänge 2022–2026.
% Aufgaben (Bank): 1 funktionsklassen-und-eigenschaften-e5-k1-s2-v2 · 2 grenzwerte-und-verhalten-im-
% unendlichen-e2-k1-s2-v1 · 3 funktionsklassen-und-eigenschaften-e4-k1-s2-v1 · 4 funktionsklassen-und-
% eigenschaften-e2-k1-s7-v1 · 5 kurvenuntersuchung-e1-k1-s4-v1 · 6 kurvenuntersuchung-e2-k2-s2-v1 ·
% 7 kurvenuntersuchung-e3-k2-s2-v1 · 8 eigen · 9 kurvenuntersuchung-e5-k1-s8-v2 · 10 extremalprobleme-
% e2-k1-s6-v1 · 11 gleichungen-loesen-e4-k1-s1-v1 · 12 rekonstruktion-von-funktionsgleichungen-e2-k1-s1-v1
% (gekürzt auf die Gleichungen).
\AtBeginDocument{\pgfplotsset{pfklein/.style={axis lines=middle,xlabel={$x$},ylabel={$y$},tick label style={font=\scriptsize},samples=120,clip=true,every axis plot/.append style={line width=0.9pt},/pgf/number format/use comma,grid=both,grid style={mbgitter}}}}
\newcommand{\kennung}{K7M}
\newcommand{\ueteil}[3]{\par\addvspace{9pt}\Needspace*{7\baselineskip}\noindent
  {\bfseries Teil #1 \textperiodcentered{} #2}\ifblank{#3}{}{\hspace{0.7em}{\scriptsize\color{mbgrau}Rand}}\par\nobreak\vspace{1pt}}
\newcommand{\ueantwort}{\par\nobreak\vspace{1pt}\noindent{\small\color{mbgrau}Antwort:}\ \rechenplatz{1}}

\begin{document}
\pfheftstil
\fancyfoot[C]{\parbox[t]{\textwidth}{\mbfhbox\par\vspace{1.5mm}%
  {\footnotesize\color{mbgrau}\kennung\hfill\thepage}}}
\pfheftkopf{Kurvenuntersuchung}{Abi GK \textperiodcentered{} Übersicht}
\noindent Löse, was du kannst. Wo du hängst, beginnt dein Teil.

\ueteil{1}{Graphen verschieben, spiegeln, strecken}{}
$f(x) = e^x$ und $g(x) = e^{-x}$. Wie geht der Graph von $g$ aus dem Graphen von $f$ hervor?
\ueantwort
\fusshilfe{\mbox{1:~Spiegelung an der $y$-Achse}}

\ueteil{2}{Verhalten für große und kleine $x$}{}
$f(x) = (4 - x) \cdot e^{x}$. Wohin gehen die Funktionswerte für $x \to +\infty$ und für $x \to -\infty$?
\ueantwort
\fusshilfe{\mbox{2:~$x \to +\infty$: $f(x) \to -\infty$; $x \to -\infty$: $f(x) \to 0$ (von oben)}}

\ueteil{3}{Symmetrie}{rand}
$f(x) = x^2 \cdot e^{-x^2}$. Ist der Graph symmetrisch? Begründe über $f(-x)$.
\ueantwort
\fusshilfe{\mbox{3:~$f(-x) = f(x)$: achsensymmetrisch zur $y$-Achse}}

\ueteil{4}{Nullstellen und Schnittpunkte mit den Achsen}{}
$f(x) = (x^2 - 4) \cdot e^{x}$. Berechne die Nullstellen.
\ueantwort
\fusshilfe{\mbox{4:~$e^x > 0$, also $x = -2$, $x = 2$}}

\ueteil{5}{Wo steigt der Graph, wo fällt er?}{}
Eine ganzrationale Funktion $f$ dritten Grades hat genau die Extrempunkte $H(-2 \mid 5)$ und $T(1 \mid -3)$. Gib ohne Rechnung an, wo $f$ steigt und wo $f$ fällt.
\ueantwort
\fusshilfe{\mbox{5:~steigt auf $]-\infty; -2]$, fällt auf $[-2; 1]$, steigt auf $[1; \infty[$}}

\ueteil{6}{Hoch- und Tiefpunkte}{}
Zeige, dass der Graph von $f(x) = x^3 - 27x + 1$ an der Stelle $x = -3$ einen Hochpunkt hat.
\ueantwort
\fusshilfe{\mbox{6:~$f'(-3) = 27 - 27 = 0$, $f''(-3) = -18 < 0$}}

\ueteil{7}{Wendepunkte}{}
Bestimme den Wendepunkt des Graphen von $f(x) = x^3 - 3x^2 + 5$.
\ueantwort
\fusshilfe{\mbox{7:~$W(1 \mid 3)$}}

\ueteil{8}{Den Graphen skizzieren}{rand}
\begin{minipage}[t]{0.56\linewidth}\vspace{0pt}\raggedright
Der Graph von $f$ hat den Tiefpunkt $T(-1 \mid -2)$ und den Hochpunkt $H(1 \mid 2)$, sonst keine Extrempunkte. Für $x \to +\infty$ gilt $f(x) \to -\infty$. Skizziere den Graphen.
\end{minipage}\hfill
\begin{minipage}[t]{0.40\linewidth}\vspace{0pt}\raggedleft
\begin{tikzpicture}\begin{axis}[pfklein,xmin=-2.5,xmax=2.5,ymin=-3,ymax=3,width=5.4cm,height=4.4cm,xtick={-2,...,2},ytick={-2,...,2}]
\addplot[draw=none,domain=-2:2] {0};
\end{axis}\end{tikzpicture}
\end{minipage}
\fusshilfe{\mbox{8:~von links oben fallend bis $T$, steigend bis $H$, dann fallend nach unten}}

\ueteil{9}{Maße aus dem Graphen}{}
Ein Brückenbogen wird durch $f(x) = -0{,}1x^2 + 2{,}5$ beschrieben, oberhalb der $x$-Achse (1 LE = 2 m). Wie breit und wie hoch ist der Bogen in Metern?
\ueantwort
\fusshilfe{\mbox{9:~Nullstellen $\pm 5$: $20$ m breit; $f(0) = 2{,}5$: $5$ m hoch}}

\ueteil{10}{Extremalprobleme}{rand}
Der Ursprung und der Punkt $P(u \mid f(u))$ mit $u > 0$ auf dem Graphen von $f(x) = (x + 3) \cdot e^{-0{,}5x}$ sind gegenüberliegende Ecken eines achsenparallelen Rechtecks. Stelle die Zielfunktion für den Flächeninhalt in Abhängigkeit von $u$ auf.
\ueantwort
\fusshilfe{\mbox{10:~$A(u) = u \cdot (u + 3) \cdot e^{-0{,}5u}$, $u > 0$}}

\ueteil{11}{Gleichungen und Ungleichungen lösen}{}
Für welche $x$ gilt $(x - 1) \cdot (x + 2)^2 > 0$?
\ueantwort
\fusshilfe{\mbox{11:~$x > 1$}}

\ueteil{12}{Funktionsgleichung aus Bedingungen}{rand}
Der Graph einer quadratischen Funktion $f(x) = ax^2 + bx + c$ verläuft durch $A(0 \mid -1)$ und hat den Tiefpunkt $T(1 \mid -3)$. Stelle die drei Gleichungen für $a$, $b$ und $c$ auf.
\ueantwort
\fusshilfe{\mbox{12:~$c = -1$; $a + b + c = -3$; $2a + b = 0$}}
\end{document}
